\documentclass[12pt]{article}
\usepackage{graphicx}
\usepackage{amsmath}
\usepackage{siunitx}
\usepackage{booktabs}
\usepackage{hyperref}

\title{Seagrass Meadows as pH Refugia from Ocean Acidification}
\author{Seagrass pH Refugia Project Team}
\date{\today}

\begin{document}

\maketitle

\begin{abstract}
This report analyzes seagrass meadows as potential refugia from ocean acidification
by examining pH, dissolved inorganic carbon (DIC), and total alkalinity (TA) data
from multiple sources including Copernicus Marine Service, GLODAP, SOCAT, and
PyCO2SYS calculations. The study identifies seagrass meadows that may provide
localized relief from ocean acidification effects through biological carbonate
chemistry modification.
\end{abstract}

\section{Introduction}
Ocean acidification poses a significant threat to marine ecosystems worldwide.
Seagrass meadows, through their photosynthetic activity, can modify local
carbonate chemistry and potentially create refugia from low pH conditions.

\section{Data Sources}
\begin{itemize}
    \item Copernicus Marine Service satellite and model data
    \item GLODAP (Global Ocean Data Analysis Project) biogeochemical data
    \item SOCAT (Surface Ocean CO2 Atlas) measurements
    \item PyCO2SYS calculations for carbonate system parameters
\end{itemize}

\section{Methods}
\subsection{Data Processing}
Data were collected, quality-controlled, and processed using custom Python scripts.
PyCO2SYS was used to calculate carbonate chemistry parameters from measured
variables.

\subsection{Analysis}
Seagrass versus open water comparisons were conducted to assess the magnitude
of pH modification by seagrass meadows. Monthly anomalies and time series
analyses were performed to identify seasonal patterns.

\section{Results}
Key findings include:
\begin{itemize}
    \item Seagrass meadows exhibit elevated pH relative to adjacent open waters
    \item The pH difference shows seasonal variability, with maximum effects
    during periods of high productivity
    \item DIC and TA show complementary patterns consistent with photosynthetic
    carbon uptake
\end{itemize}

\begin{figure}[htbp]
    \centering
    \includegraphics[width=\textwidth]{figures/delta_ph_timeseries.png}
    \caption{Seagrass--Open Water \textit{Δ}pH: 2019--20}
    \label{fig:delta_ph}
\end{figure}

\section{Discussion}
The observed pH elevations in seagrass meadows suggest their potential to act as
local refugia from ocean acidification. The magnitude and variability of this
effect have implications for marine conservation and climate change adaptation
strategies.

\section{Conclusion}
Seagrass meadows demonstrate measurable capacity to modify local carbonate
chemistry, providing potential buffering against ocean acidification effects.
Further research is needed to understand the spatial and temporal scales of
this refugia effect.

\section*{Acknowledgments}
This project utilized data from Copernicus Marine Service, GLODAP, SOCAT, and
involved PyCO2SYS calculations for carbonate system parameters.

\bibliography{references}
\bibliographystyle{unsrt}

\end{document}